%!TeX root=MemoriaTFG.tex

\chapter{Introducció}

\section{Contexto}

Como participante activo que he sido en \ac{LVUIB} durante los años que se ha celebrado me he dado cuenta de lo costoso que es encontrar un sistema a nivel de aplicación que permita organizar y gestionar una liga deportiva. Suelen ser necesarias muchas funcionalidades avanzadas y personalizadas que llevan a optar por soluciones que acaban necesitando una inversión monetaria, pero que aún así no acaban de convencer a los participantes y usuarios de dicho sistema. 

Esta situación es una descripción de lo ocurrido con la \ac{LVUIB} del año 2025-2026. La experiencia de los alumnos participantes con dicha aplicación resultó en un descontento general debido a: el uso de múltiples aplicaciones para los árbitros, la incorrecta persistencia y tardía sincronización de los resultados de los partidos y la falta de funcionalidades que obligaba a mantener un flujo de gestión híbrido entre la aplicación y la propia \ac{UIB} entre otras razones.

Cabe remarcar que la \ac{OUSIS} es el organismo encargado de fomentar y gestionar los eventos deportivos y relacionados con la salud que se celebran dentro del entorno de la \ac{UIB}, como pueden ser las ligas universitarias de fútbol y pádel. Esta gestión no estuvo tan presente durante el primer año que se celebró la \ac{LVUIB}, ya que fue un evento fomentado y gestionado por una asociación de alumnos creada al uso para dicho evento. En años posteriores, a medida que la \ac{LVUIB} fue creciendo en popularidad y participantes, la \ac{OUSIS} decidió involucrarse de lleno.

Debido a esta situación llegué a la conclusión de que era posible desarrollar un sistema que permitiera una gestión centralizada evitando los problemas comentados y que, además, permitiera utilizar la aplicación como el único punto necesario de gestión.

\section{Motivación y objetivos}

La motivación principal de este \ac{TFG} ha sido la de crear un sistema gestor de ligas deportivas que aporte una propuesta de valor con respecto a las soluciones que existen en el mercado y de las que he sido usuario a nivel personal.

\begin{enumerate}[label=\textbf{O\arabic*}, leftmargin=1.5cm, sep=1ex]
    \item \textbf{Desarrollar la aplicación local:} Diseñar e implementar el entorno de desarrollo utilizando herramientas eficientes en el sistema operativo.
    \item \textbf{Optimizar el rendimiento:} Reducir los tiempos de compilación del sistema mediante la automatización de tareas secundarias.
    \item \textbf{Evaluar la experiencia de usuario:} Realizar pruebas periódicas para garantizar que la interfaz cumple con los estándares formales.
\end{enumerate}

\section{Alternativas actuales}

Aunque actualmente podemos encontrar miles de alternativas que aportan las funcionalidades básicas y genéricas necesarias para una correcta gestión de ligas deportivas, prefiero centrarme en las aplicaciones utilizadas en las temporadas anteriores de la \ac{LVUIB}. Dichas aplicaciones son: Winner y Clupik

% TODO: Añadir referencias a Winner y Clupik

Winner se utilizó en el primer y segundo año que se celebró la \ac{LVUIB}. Es una aplicación que se puede encontrar en cualquier tienda de aplicaciones móviles. Dado que la responsabilidad de gestión cayó principalmente en la asociación de alumnos anteriormente comentada, se optó por una herremienta gratuita que permitía funcionalidades básicas: gestión de equipos, gestión de partidos y visualización de tablas clasificación. En caso de querer mas funcionalidades, se debía hacer una pequeña inversión no superior a 10€: una opción asequible y por la que se optó en el segundo año.

En cambio, Clupik se ha estado utilizando durante el tercer año (coincidiendo con el año académico 2025-2026) y ha representado un cambio importante respecto a la primera alternativa, ya que la propia \ac{OUSIS} fue la responsable de adquirir la herramienta. Esta herramienta se presenta como un \ac{SaaS} por el que la \ac{OUSIS} ha tenido que hacer una inversión entorno a los 600€. Este coste se debe a que Clupik es una solución profesional, personalizada y enfocada a organismos que necesitan sistemas gestores de ligas deportivas.
Podría parecer que esta herramienta es la ideal para el caso de la \ac{LVUIB}, pero todo lo contrario. Esta alternativa ha presentado muchos problemas: tanto por parte de los jugadores como de los árbitros y administradores de la liga.

